\documentclass[11pt]{article}
\usepackage[T1]{fontenc}
\usepackage[utf8]{inputenc}
\usepackage{lmodern}
\usepackage[margin=1in]{geometry}
\usepackage{amsmath,amssymb,amsthm}
\usepackage{microtype}
\usepackage[colorlinks=true,linkcolor=blue,citecolor=blue,urlcolor=blue]{hyperref}
\hypersetup{pdftitle={A Narayana factorization for a truncated two-row grid polytope},
  pdfauthor={Xiaojian Zeng},pdfkeywords={Narayana polynomial, preorder polytope, lattice point, rank-one commutator}}
\newtheorem{theorem}{Theorem}
\newtheorem{lemma}[theorem]{Lemma}
\DeclareMathOperator{\CT}{CT}
\DeclareMathOperator{\tr}{tr}
\DeclareMathOperator{\supp}{supp}
\newcommand{\A}{\mathsf A}
\newcommand{\B}{\mathsf B}
\newcommand{\T}{\mathsf T}
\newcommand{\Q}{\mathbb Q}
\newcommand{\Z}{\mathbb Z}
\title{A Narayana factorization for a truncated two-row grid polytope}
\author{Xiaojian Zeng\\[3pt]
\small PBC School of Finance, Tsinghua University, 100190\\
\small \href{mailto:xiaojianzeng@foxmail.com}{xiaojianzeng@foxmail.com}}
\date{}

\begin{document}
\maketitle
\begin{abstract}
Remove the maximum element from the product of an $n$-element chain and
a two-element chain. For the associated polytope, defined by nonnegativity
and an upper bound of $|I|$ on the sum over each lower ideal $I$, we prove
that the lattice-point support enumerator is the product of the type~A
Narayana polynomial of index $n-1$ and the type~B Narayana polynomial of
index $n$. This establishes the factorization proposed by Athanasiadis
and Chapoton. A transfer with two slack variables becomes a pair of
operators with a rank-one commutator. An explicit finite boundary trace,
followed by formal Lagrange inversion and telescoping, evaluates its
terminal moment for every $n\geq1$.
\end{abstract}

\section{The support identity}
Let $[n]=\{1,\ldots,n\}$ and let
\[
 P_n=([n]\times[2])\setminus\{(n,2)\},
 \qquad (i,r)\leq(j,s)\ \Longleftrightarrow\ i\leq j\text{ and }r\leq s.
\]
A lower ideal contains every predecessor of each of its elements. Set
\[
 Q_n=\left\{x\in\mathbb R_{\geq0}^{P_n}:
          \sum_{v\in I}x_v\leq |I|
          \text{ for every lower ideal }I\subseteq P_n\right\},
 \qquad
 H_n(t)=\sum_{x\in Q_n\cap\Z^{P_n}}t^{|\supp(x)|},
\]
where $\supp(x)=\{v:x_v\ne0\}$. The full ideal bounds the total coordinate
sum by $2n-1$, so this is a finite sum. The empty ideal adds no condition.
For $r\geq0$, use the normalizations
\begin{equation}\label{eq:narayana}
 A_r(t)=\sum_{j=0}^r\frac{1}{j+1}\binom rj\binom{r+1}j t^j,
 \qquad B_r(t)=\sum_{j=0}^r\binom rj^2t^j.
\end{equation}

\begin{theorem}\label{thm:main}
For every integer $n\geq1$,
\begin{equation}\label{eq:main}
 H_n(t)=A_{n-1}(t)B_n(t).
\end{equation}
\end{theorem}

This is the maximum-deleted grid identity proposed in the final paragraph
of \cite[Section~9.3]{AC26}. The polynomial denoted by $h$ there counts
nonzero coordinates, as defined in \cite[Section~5]{AC26}; the present
statement concerns this support statistic. The two factors already occur
separately in composition polytopes: a total-sum bound gives $B_r$, while
all chain-prefix bounds give $A_r$ \cite[Examples~2.3--2.4]{Ath25}.
Those evaluations do not by themselves separate the coupled grid
inequalities. We evaluate the coupled transfer directly, without requiring
a support-preserving product bijection.

All series below are formal. We first introduce an independent parameter
$c$, work over $\Q(c,t)$ when division by $t$ is convenient, and recover
polynomial identities before specializing $c=1+t$.

\section{Two slack variables and the transfer}
Write $a_i=x_{(i,1)}$ and $b_j=x_{(j,2)}$. An ideal is exactly a pair of
row prefixes of lengths $(i,j)$, where
$0\leq i\leq n$ and $0\leq j\leq\min(i,n-1)$.
Indeed, each row must be a prefix, and inclusion of $(j,2)$ forces
inclusion of $(j,1)$; these conditions also suffice. Define
\[
 U_i=i-\sum_{k=1}^ia_k,\qquad V_j=j-\sum_{k=1}^jb_k,
 \qquad U_0=V_0=0.
\]
The ideal inequalities are therefore precisely $U_i+V_j\geq0$ over the
indicated pairs. After $m$ complete columns, record
\[
 M_m=\min(V_0,\ldots,V_m),\qquad z=U_m+M_m,\qquad s=V_m-M_m.
\]
Both $z$ and $s$ are nonnegative. Appending the next two coordinates
$a,b$ is legal exactly when
\begin{equation}\label{eq:bounds}
 0\leq a\leq z+1,\qquad 0\leq b\leq z+s+2-a.
\end{equation}
The first bound imposes the new lower prefix against every old upper
prefix. The second imposes the pair of new prefixes, which is the only
remaining new ideal. If $\delta=\max(0,b-s-1)$, the new minimum is
$M_m-\delta$, and the new state is
\begin{equation}\label{eq:update}
 (z+1-a-\delta,\ \max(0,s+1-b)).
\end{equation}
Thus the two recorded slacks suffice for every future constraint.

Give each positive coordinate weight $t$ and each zero coordinate weight
$1$. If $F_m(x,y)$ is the state polynomial after $m$ complete columns,
the two half-step operators are
\begin{align}
 \alpha(x^zy^s)&=\left(x^{z+1}+t\sum_{r=0}^z x^r\right)y^s,
 \label{eq:alpha}\\
 \beta(x^ry^s)&=x^r\left(y^{s+1}+t\sum_{h=0}^s y^h\right)
                    +t\sum_{h=0}^{r-1}x^h.\label{eq:beta}
\end{align}
The last sum is empty when $r=0$. In \eqref{eq:beta}, the three groups
correspond to $b=0$, $1\leq b\leq s+1$, and
$s+2\leq b\leq s+1+r$, respectively. Starting at $(0,0)$ gives
$F_m=(\beta\alpha)^m1$. After $n-1$ complete columns only $a_n$ remains,
so
\begin{equation}\label{eq:terminal}
 H_n(t)=\bigl(\alpha(\beta\alpha)^{n-1}1\bigr)(1,1).
\end{equation}

Make the change of variables $x=1+u$, $y=1+v$. On polynomials in $u,v$,
let $\mathsf U_i$ multiply by the $i$th variable, let $\mathsf D_i$
lower its exponent by one and kill its constant term, and let
$\mathsf E_i$ set that variable to zero. Then
\[
 \mathsf D_i\mathsf U_i=1,\qquad
 \mathsf U_i\mathsf D_i=1-\mathsf E_i.
\]
Introduce
\begin{equation}\label{eq:operators}
 \A=\mathsf U_1+c+t\mathsf D_1,\qquad
 \B=\mathsf U_2+c+t\mathsf D_2+t\mathsf D_1\mathsf E_2,
 \qquad \T=\B\A.
\end{equation}
At $c=1+t$, these are the transformed $\alpha,\beta$. For example,
the transformed first half-step sends $f$ to
$(1+u+t)f+t(f-f(0,v))/u$. The second sends it to
\[
 (1+v+t)f+t\frac{f-f(u,0)}{v}
                    +t\frac{f(u,0)-f(0,0)}{u},
\]
which verifies every term in \eqref{eq:operators}.
Let $e=1$, let $\epsilon f=[u^0v^0]f$, and put
$E=\mathsf E_1\mathsf E_2=e\epsilon$. Equation~\eqref{eq:terminal} becomes
\begin{equation}\label{eq:moment}
 h_n:=\epsilon\A\T^{n-1}e,\qquad H_n(t)=h_n\big|_{c=1+t}.
\end{equation}
With $[X,Y]=XY-YX$, the only noncommuting summands in $\B,\A$ give
\begin{equation}\label{eq:commutator}
 [\B,\A]=t[\mathsf D_1,\mathsf U_1]\mathsf E_2=tE,
 \qquad [\A,\T]=-tE\A.
\end{equation}

\section{The finite boundary trace}
We now keep $c,t$ independent. In the basis $u^iv^j$, $i,j\geq0$, the
operators have bounded displacement. A fixed word containing $E$ has
only finitely many nonzero matrix entries: both its input and output
indices must be within a bounded distance of the origin. Its trace and
the cyclic permutations used below are consequently finite sums.
Expanding \eqref{eq:commutator} gives
\begin{equation}\label{eq:trace}
 \tr[\A,\T^n]
 =-t\sum_{r=0}^{n-1}\tr(\T^rE\A\T^{n-1-r})=-nt h_n.
\end{equation}
Here and below, we take the trace of the displayed finite-support
operator, without assigning separate traces to its infinite summands.

Set
\[
 S(u)=u+c+t/u,\qquad J=\mathsf U+c+t\mathsf D
       \quad\hbox{on polynomials in }v,
 \qquad B_u=J+(t/u)E_v,
\]
where $E_vf=f(0)$. Write $\CT_u$ for the coefficient of $u^0$.

\begin{lemma}\label{lem:boundary}
For every $n\geq1$, with
$R_n(y)=\tr((J+yE_v)^n-J^n)$,
\begin{equation}\label{eq:boundary}
 nt h_n=\CT_u (u-t/u)S(u)^nR_n(t/u).
\end{equation}
\end{lemma}
\begin{proof}
Put $K=\T^n$, with the first matrix index denoting the output. For each
fixed $j$, cancellation of adjacent horizontal terms gives
\begin{equation}\label{eq:flux}
 \sum_{i=0}^N[\A,K]_{(i,j),(i,j)}
       =tK_{(N+1,j),(N,j)}-K_{(N,j),(N+1,j)}.
\end{equation}
There is no lower-boundary contribution because the unilateral lowering
operator kills exponent zero. Choose $N\geq2n+2$. Each of the $2n$
half-steps in $K$ moves the horizontal index by at most one. Thus paths
contributing to the two entries on the right stay away from horizontal
index zero, and their horizontal shifts may be treated as bilateral.
The Laurent symbol of $K$ there is $S(u)^nB_u^n$. The two entries are its
$u^1$ and $u^{-1}$ coefficients, respectively, on vertical diagonal $j$.
The right side of \eqref{eq:flux} is therefore
\[
 -\CT_u(u-t/u)S(u)^n(B_u^n)_{j,j}.
\]
Sum over $0\leq j\leq L$, where $L\geq n+2$. This rectangle contains
the diagonal support of \eqref{eq:trace}: its terms propagate at most
$n$ vertical steps and $2n+1$ horizontal steps from the origin.
For each $j$ we may replace $(B_u^n)_{j,j}$ by
$(B_u^n-J^n)_{j,j}$, since
\[
 \CT_u(u-t/u)S(u)^n
   =\frac{1}{n+1}\CT_u\,u\frac{d}{du}S(u)^{n+1}=0.
\]
Every word in $B_u^n-J^n$ contains $E_v$, so its diagonal entries vanish
for $j>n$. The resulting vertical sum is a finite trace and is independent
of $L$. Comparing with \eqref{eq:trace} proves \eqref{eq:boundary},
including its sign.
\end{proof}

\section{Formal extraction and the two factors}
Let $\epsilon_v$ extract the constant term in $v$, and define
\[
 M(q)=\epsilon_v(1-qJ)^{-1}1.
\]
Its coefficients enumerate nonnegative Motzkin excursions with up-weight
$1$, level-weight $c$, and down-weight $t$. This is the constant-weight
case of the weighted-path framework of \cite[Theorem~1]{Flajolet80}.
Decomposing at the first step and, for an up-step, its first matching
return gives directly
\begin{equation}\label{eq:motzkin}
 M=1+cqM+tq^2M^2.
\end{equation}
Consequently $V=tqM$ is the unique series with zero constant term satisfying
\begin{equation}\label{eq:lagrange}
 V=q(V^2+cV+t).
\end{equation}
The rank-one perturbation has the logarithmic generating series
\begin{equation}\label{eq:log}
 \sum_{n\geq1}\frac{R_n(y)}{n}q^n=-\log(1-yqM(q)).
\end{equation}
To verify this without infinite traces, differentiate the left side in
$y$. Cyclicity of the finite-support terms yields
$q\epsilon_v(1-q(J+yE_v))^{-1}1$. Expansion by the number of insertions
of $E_v=1\epsilon_v$ evaluates this as $qM/(1-yqM)$. Integrating
coefficientwise and using the value zero at $y=0$ proves \eqref{eq:log}.

Substituting $y=t/u$ into \eqref{eq:log} gives
$R_n(t/u)=-n[q^n]\log(1-V/u)$. Lemma~\ref{lem:boundary} and formal
constant-term integration by parts now imply
\begin{align}
 t h_n
 &=-[q^n]\CT_u(u-t/u)S(u)^n\log(1-V/u)\notag\\
 &=\frac{1}{n+1}[q^n]\CT_u S(u)^{n+1}\frac{V}{u-V}.
 \label{eq:extract}
\end{align}
Indeed, $u(d/du)S^{n+1}=(n+1)(u-t/u)S^n$ and
$u(d/du)\log(1-V/u)=V/(u-V)$. Expand the last fraction in powers of
$V/u$. Since $V$ has zero constant term, each $q$ coefficient involves
only finitely many Laurent monomials, justifying these operations.

For every integer $k$, put $C_{m,k}=[u^k]S(u)^m$.
Formal Lagrange inversion \cite[Theorem~2.1.1]{Gessel16}, applied to
\eqref{eq:lagrange}, gives for $1\leq k\leq n$
\begin{equation}\label{eq:lagrange-coeff}
 [q^n]V^k=\frac{k}{n}[w^{n-k}](w^2+cw+t)^n
         =\frac{k}{n}C_{n,-k}=\frac{k}{n}t^kC_{n,k}.
\end{equation}
The last equality follows from $S(t/u)=S(u)$; coefficients vanish for
$k>n$. For completeness, the inversion formula here also follows by
writing $\phi(w)=w^2+cw+t$ and changing variables
$q=w/\phi(w)$ in the formal residue for $[q^n]V^k$.
The residue is
\[
 \mathop{\rm Res}_w
 w^{k-n-1}\phi(w)^n\left(1-\frac{w\phi'(w)}{\phi(w)}\right).
\]
The zero residue of $(w^{k-n}\phi(w)^n)'$ makes this
$(k/n)[w^{n-k}]\phi(w)^n$. The substitution is valid over
$\Q(c,t)$ because $\phi(0)=t$ is invertible there.

Expanding $V/(u-V)=\sum_{k\geq1}V^ku^{-k}$ in \eqref{eq:extract}
therefore yields
\begin{equation}\label{eq:sum}
 h_n=\frac{1}{n(n+1)t}
       \sum_{k=1}^n k t^k C_{n+1,k}C_{n,k}.
\end{equation}
Differentiating $S^{n+1}$ gives
\[
 kC_{n+1,k}=(n+1)(C_{n,k-1}-tC_{n,k+1}).
\]
Thus the numerator of \eqref{eq:sum} telescopes to
\begin{align*}
 &(n+1)\sum_{k\geq1}
 \left(t^k C_{n,k-1}C_{n,k}-t^{k+1}C_{n,k}C_{n,k+1}\right)\\
 &=(n+1)tC_{n,0}C_{n,1}.
\end{align*}
All sums are finite because $C_{n,k}=0$ for $k>n$. We have proved
\begin{equation}\label{eq:generic}
 \epsilon\A(\B\A)^{n-1}1=\frac{1}{n}C_{n,0}C_{n,1}.
\end{equation}
Both sides are polynomials in $c,t$, so the identity holds also at $t=0$.

\begin{proof}[Completion of the proof of Theorem~\ref{thm:main}]
Specialize $c=1+t$. Then $S(u)=(1+u)(1+t/u)$, and binomial expansion gives
\begin{align*}
 C_{n,0}&=\sum_{j=0}^n\binom nj^2t^j=B_n(t),\\
 \frac{C_{n,1}}{n}
 &=\sum_{j=0}^{n-1}\frac{1}{n}\binom n{j+1}\binom nj t^j
 =\sum_{j=0}^{n-1}\frac{1}{j+1}\binom{n-1}j\binom nj t^j
 =A_{n-1}(t).
\end{align*}
Equations~\eqref{eq:moment} and \eqref{eq:generic} prove \eqref{eq:main}.
For $n=1$ this gives $1+t$, agreeing with $Q_1=[0,1]$. At $t=0$ both
sides are $1$, and no exceptional specialization is excluded.
\end{proof}

\section{Exact checks and computational disclosure}
The proof is accompanied by two separately written Python verifiers
using integer arithmetic. Reconstructing the finite objects in a separate
exact checker is a verification practice also used in
\cite[Section~6]{ZengEtAl26}; here the checker reconstructs ideals from
the coordinatewise order, independently of the transfer implementation.
For $n=1,\ldots,6$, it tests every nonnegative integer vector of total
at most $2n-1$ against every ideal. The ambient counts are
$2,20,252,3432,48620,705432$, respectively, and all six complete support
polynomials agree with \eqref{eq:main}.

Further checks test the commutator on all $u^iv^j$ with $0\leq i,j\leq7$,
the moment and coefficient-sum identities through $n=16$, and the finite
boundary identity through $n=8$, at each parameter pair
\[
 (c,t)\in\{(0,0),(2,3),(0,1),(-3,2),(4,-3),(1,0),(-2,-2)\}.
\]
These are exact finite checks of the formulas, including zero and negative
weights. They do not extend the quantified ranges of the computations;
the all-$n$ result follows from the formal argument above.

OpenAI Codex assisted with proof development, literature searches, Python
code, a separate model-context review of the argument, and manuscript
preparation. The second checker was written in that separate context
without importing the first implementation. This separation describes
the computation and review procedure; it is not external peer review or
a claim of statistical independence. No proof assistant was used. The
accompanying source archive records the check commands and their exact
coverage.

\bibliographystyle{plain}
\bibliography{references}
\end{document}
